\section*{Chapitre \ref{ch:g12:catalysis} --- La catalyse}

\begin{solution}{exo:g12:catalysis:1}
\omterm{def:g7:chemical-reactions:reactant}{Réactif}~: \ce{H2O2}. Produits~: \ce{H2O} et \ce{O2}. \omterm{def:g12:catalysis:catalyst}{Catalyseur}~:
\ce{I-} (consommé lors de l'étape 1, restitué lors de l'étape 2).
Intermédiaire~: \ce{IO-} (formé lors de l'étape 1, consommé lors de
l'étape 2).
\end{solution}

\begin{solution}{exo:g12:catalysis:2}
Homogène~; hétérogène~; hétérogène~; homogène.
\end{solution}

\begin{solution}{exo:g12:catalysis:3}
Le chemin non catalysé. Les niveaux de départ et d'arrivée sont les
mêmes~: les mêmes produits se forment, avec la même énergie libérée~;
seul le chemin diffère.
\end{solution}

\begin{solution}{exo:g12:catalysis:4}
Il est consommé lors d'une étape et restitué lors d'une autre~: il
apparaît des deux côtés de la somme des étapes, et s'élimine.
\end{solution}

\begin{solution}{exo:g12:catalysis:5}
La réaction a lieu à la surface~: plus la surface est grande, plus il y
a de \omterm{def:g7:atoms-and-molecules:molecule}{molécules} qui peuvent réagir à la fois, pour une même masse de
\omterm{def:g12:catalysis:catalyst}{catalyseur} (souvent coûteux).
\end{solution}

\begin{solution}{exo:g12:catalysis:6}
Somme~: \ce{2Fe^{3+} + 2H2O2 + 2Fe^{2+} + 2H+ -> 2Fe^{2+} + O2 + 2H+ +
2Fe^{3+} + 2H2O}. Les \omterm{def:g9:ions:ion}{ions} fer et les \omterm{def:g9:ions:ion}{ions} hydrogène apparaissent des
deux côtés et s'éliminent~: \ce{2H2O2 -> 2H2O + O2}.
\end{solution}

\begin{solution}{exo:g12:catalysis:7}
La moitié de \qty{72}{mL}, soit \qty{36}{mL}, est atteinte au bout
d'environ \qty{35}{min} sans \omterm{def:g12:catalysis:catalyst}{catalyseur} et de \qty{3.5}{min} avec. Le
\omterm{def:g12:catalysis:catalyst}{catalyseur} raccourcit cette durée d'un facteur dix.
\end{solution}

\begin{solution}{exo:g12:catalysis:8}
La réaction a lieu à la surface du platine~; le plomb la recouvre et y
reste, si bien que les gaz ne peuvent plus atteindre le \omterm{def:g9:periodic-table-first-look:metal}{métal}. Le
\omterm{def:g12:catalysis:catalyst}{catalyseur} est «~empoisonné~»~: présent, mais inutile.
\end{solution}

\begin{solution}{exo:g12:catalysis:9}
Jusque vers \qty{37}{\celsius}, la réaction s'accélère avec la
température, comme toute réaction. Au-delà, l'\omterm{def:g12:catalysis:enzyme}{enzyme}, une protéine, perd
sa forme et avec elle son activité~; à \qty{80}{\celsius}, elle est
détruite.
\end{solution}

\begin{solution}{exo:g12:catalysis:10}
\ce{2CO + 2NO -> 2CO2 + N2}~: C 2 = 2, O 4 = 4, N 2 = 2. Le monoxyde de
carbone est oxydé par le monoxyde d'azote, qui est réduit~: chaque
polluant élimine l'autre.
\end{solution}

\begin{solution}{exo:g12:catalysis:11}
$n(\ce{O2}) = 0.072 / 24.1 = \qty{2.99e-3}{mol}$~;
$n(\ce{H2O2}) = 2 \times \num{2.99e-3} = \qty{5.98e-3}{mol}$ dans
\qty{10.0}{mL}~: $c = \qty{0.598}{mol/L}$.
\end{solution}

\begin{solution}{exo:g12:catalysis:12}
\ce{C2H5OH -> C2H4 + H2O}, une \omterm{def:g12:curly-arrows:categories}{élimination}~; \ce{C2H5OH -> CH3CHO + H2}.
Un \omterm{def:g12:catalysis:selective}{catalyseur sélectif} accélère l'une de plusieurs réactions possibles
beaucoup plus que les autres, et décide ainsi du produit.
\end{solution}

\begin{solution}{exo:g12:catalysis:13}
Le \omterm{def:g12:catalysis:catalyst}{catalyseur} ne change pas l'\omterm{def:g10:reaction-progress-table:system}{état final}~: les courbes avec et sans
\omterm{def:g12:catalysis:catalyst}{catalyseur} aboutissent au même volume de dioxygène. On obtient la même
quantité de produit, seulement plus tôt.
\end{solution}

\begin{solution}{exo:g12:catalysis:14}
$\num{6.02e23} / \num{5e6} = \qty{1.2e17}{s}$, quelque quatre milliards
d'années. Pour le faire en une seconde~: \num{1.2e17} \omterm{def:g7:atoms-and-molecules:molecule}{molécules}
d'\omterm{def:g12:catalysis:enzyme}{enzyme}, soit
\qty{2e-7}{mol}.
\end{solution}

\begin{solution}{exo:g12:catalysis:15}
Ammoniac \ce{NH3}~: $150 \times 17.0 / 14.0 = \qty{182}{Mt}$. \omterm{def:g7:atoms-and-molecules:atom}{Atomes}
d'azote~: $\num{1.50e14} / 14.0 = \qty{1.07e13}{mol}$, soit
\qty{5.36e12}{mol} de \ce{N2}.
\end{solution}

\begin{solution}{pb:g12:catalysis:1}
\textbf{1.} \ce{2H2O2 -> 2H2O + O2}.

\textbf{2.} Oui~: le peroxyde d'hydrogène est l'\omterm{def:g11:redox:oxidant}{oxydant} du couple
\ce{H2O2}/\ce{H2O} et le \omterm{def:g11:redox:oxidant}{réducteur} du couple \ce{O2}/\ce{H2O2}~; il
s'oxyde et se réduit lui-même.

\textbf{3.} Sans \omterm{def:g12:catalysis:catalyst}{catalyseur}, la réaction est extrêmement lente.

\textbf{4.} $20 / 22.4 = \qty{0.893}{mol}$.

\textbf{5.} Deux \omterm{def:g10:the-mole:amount}{moles} de peroxyde d'hydrogène par \omterm{def:g10:the-mole:amount}{mole} de dioxygène~:
\qty{1.79}{mol}.

\textbf{6.} \qty{1.79}{mol/L}.

\textbf{7.} $M = 2 \times 1.0 + 2 \times 16.0 = \qty{34.0}{g/mol}$;
$C_m = 1.79 \times 34.0 = \qty{60.7}{g/L}$.

\textbf{8.} La moitié~: \qty{0.893}{mol/L}.

\textbf{9.} La catalase~: une \omterm{def:g12:catalysis:enzyme}{enzyme}, dissoute, un \omterm{def:g12:catalysis:catalyst}{catalyseur}
biologique~; les \omterm{def:g9:ions:ion}{ions} fer(III)~: \omterm{def:g12:catalysis:homogeneous}{catalyse homogène}.

\textbf{10.} \ce{2Fe^{3+} + H2O2 -> 2Fe^{2+} + O2 + 2H+} et
\ce{2Fe^{2+} + H2O2 + 2H+ -> 2Fe^{3+} + 2H2O}~; leur somme est
\ce{2H2O2 -> 2H2O + O2}.

\textbf{11.} Avec un \omterm{def:g12:catalysis:catalyst}{catalyseur}, le volume augmente beaucoup plus vite~;
les deux courbes tendent vers le même volume final.

\textbf{12.} Les \omterm{def:g9:ions:ion}{ions} fer(III) (orange) sont transformés en \omterm{def:g9:ions:ion}{ions} fer(II)
(vert pâle) lors de la première étape, puis redeviennent fer(III) lors
de la seconde~; à la fin, tout le fer est de nouveau à l'état de
fer(III).

\textbf{13.} La cuisson a détruit l'\omterm{def:g12:catalysis:enzyme}{enzyme}.

\textbf{14.} $0.100 \times 20 = \qty{2.0}{L}$.

\textbf{15.} $n(\ce{O2}) = 0.100 \times 0.893 = \qty{0.0893}{mol}$,
soit $0.0893 \times 24.1 = \qty{2.15}{L}$.

\textbf{16.} La lente décomposition libère du dioxygène, qui ferait
monter la pression dans un flacon étanche.

\textbf{17.} $1.50 / 2 = \qty{0.750}{mol}$ de dioxygène par litre, soit
$0.750 \times 22.4 = \qty{16.8}{L}$~: «~16.8 volumes~».

\textbf{18.} $(1.79 - 1.50)/1.79 \approx \qty{16}{\percent}$.

\textbf{19.} \qty{1.79}{mol/L}.
\end{solution}
